\section{Introduction}
\label{sec:introduction}

Silicon Valley Bank's annual report for 2021 told its readers what a rate rise would
do to the bank. It disclosed an economic value of equity of \$20,672m and a profile
of that value under four parallel shocks: a 200 basis point rise would remove
\$5,722m, or 27.7\%, of it \citep{SVB10K2021}. Set against the \$16,205m of Tier~1
capital reported for the same date, that is a loss of 35.3\% of Tier~1, more than
twice the 15\% threshold at which the Basel Committee's standard identifies a bank as
warranting supervisory attention \citep{BCBS2016}. The arithmetic requires no model:
both figures appear in the same annual report. The report for 2022, filed on
24~February 2023 and two weeks before
the bank was closed, contains net-interest-income sensitivity at the same four shocks
and no economic-value figure at any of them \citep{SVB10K2022}.

That pairing invites a question with a determinate answer. If an analyst at the end of
2021 had taken the disclosed profile seriously and tried to rebuild it, using
nothing but the filing itself and public curve data, would the reconstruction have
held up? And would the same framework, applied without modification to a bank that
carried comparable securities losses through 2022 and survived, have distinguished
between them?

This paper answers both. The reconstruction holds up on the asset side and the
framework does distinguish the two banks, but the exercise surfaces something that
matters more than either result: a disclosed economic value of equity is not a
measurement in the way a disclosed securities fair value is, because it nets an asset
valuation the filings constrain tightly against a deposit valuation they do not
constrain at all.

That proposition is stated here as a general one, and it should be read as a
conjecture this case motivates rather than a result it establishes. The evidence is a
single bank-year: SVB's FY2021 report is the only observation in this sample that
carries a disclosed economic-value profile at all, which is itself part of the
finding. Whether the same gap between a constrained asset side and an unconstrained
liability side characterises such disclosures generally is a question for a sample
this paper does not have.

\subsection{What is done}

A valuation model is built from the inline XBRL of four annual reports, two for SVB
and two for the Charles Schwab Corporation, together with Treasury par yield curves
and public short-rate and mortgage-rate series. Zero curves are bootstrapped and
fitted with the Svensson form \citep{NelsonSiegel1987, Svensson1994}, the short rate
follows a one-factor Hull--White process \citep{HullWhite1990}, agency pass-throughs
carry a prepayment model, and non-maturity deposits are valued as a decaying
below-market funding source in the sense of \citet{Drechsler2021}. Everything is
calibrated on 31~December 2021 information, one option-adjusted spread per
security-type and maturity bucket, and the spreads are then frozen.

The test is to price the same portfolios at 31~December 2022 under the curve that
actually obtained and compare against what the banks subsequently disclosed. SVB's
held-to-maturity book comes out at 80.36 per \$100 of carrying amount against the
83.41 disclosed, an error of $-3.66\%$. Schwab requires a correction first, because it
moved \$206,828m of original cost out of available-for-sale during 2022 and rebased
its reported basis; measured against what it actually paid, the error is $-1.75\%$
rather than the $-7.16\%$ the uncorrected comparison suggests.

The model is then asked to reproduce SVB's published economic-value profile. Holding
the deposit beta at the 0.60 SVB itself discloses, one parameter is solved so the
model returns the $-13.9\%$ reported at $+100$
basis points. At $+200$ basis points, which the model was not fitted to, it returns
$-27.61\%$ against the $-27.70\%$ disclosed. It also recovers a feature of the
profile's shape: SVB's disclosure is non-monotonic downward, adding 6.0\% of economic
value at $-100$ basis points but only 3.1\% at $-200$, and the model independently
turns in the same place, because the standard's post-shock floor absorbs most of the
second hundred basis points at a date when the entire front of the curve sat below
2\%.

\subsection{What the exercise shows}

The half-life that reproduces SVB's disclosure is 3.59~years. That single number is
the paper's subject, for two reasons.

It is not recoverable from any filing. Neither bank disclosed how long its
non-maturity deposits were assumed to remain, and the measured exposure is acutely
sensitive to the choice. A two-year half-life is no less defensible for a base that
was 95.5\% uninsured at the end of 2022 and concentrated in venture-backed
technology, and on that assumption the same balance sheet reports $-62.47\%$ at $+200$ basis points
instead of $-27.61\%$. Across the grid of deposit betas and half-lives the figure runs
from $-122\%$ to $+21\%$. The published sensitivity was the output of an assumption
about depositor loyalty that the filing did not state and no reader could have
recovered.

And it is the assumption that the comparison with Schwab isolates. Both banks carried
similar unrecognised held-to-maturity losses into the end of 2022, \$15,152m against
\$14,138m, and both breach the supervisory outlier screen in every bank-year tested.
On a run model that meets outflows from cash, then secured borrowing, then sales, only
SVB crosses a threshold: 67.5\% of uninsured deposits at the end of 2022, binding
through liquidity rather than solvency. Schwab, funded 18.8\% by uninsured money
against SVB's 95.5\%, crosses neither threshold in either year under any of 42
parameter combinations tried. What separates them on this model is the funding
structure rather than the size of the securities loss.

\subsection{Relation to existing work}

\citet{Metrick2024} gives the standard account of the failure and the panic that
followed it. Two results frame the mechanism this paper models.
\citet{Drechsler2021} establish that a deposit franchise hedges a long-duration asset
book precisely to the extent that the deposits are sticky and insensitive to rates,
which is why an assumption about depositor behaviour can dominate a measured
exposure. \citet{Jiang2024} show that mark-to-market losses translate into solvency
problems in proportion to uninsured leverage, and mark the U.S.\ banking system to
market cross-sectionally to identify banks at risk of self-fulfilling runs.
\citet{Choi2023} supply the market counterpart, finding that investors appeared to
price the risk of uninsured deposit reliance while failing to anticipate the
realisation of the implied losses, which is a partial answer to the question of what
public information already contained.

The claim that a deposit assumption is a decision rather than a measurement is not
new, and the priority belongs elsewhere. \citet{Chen2026} argues exactly this for
uninsured balances in the banking book: that the retention schedule used in an
economic-value calculation is a management choice, that its alternatives should be
priced alongside it, and that the choice moves present value and duration across
ranges wide enough to make the disclosed figure uninformative on its own. That work
is conducted on a synthetic book, with no bank data and no named institution.

What this paper adds is the instantiation. It takes a real bank that published an
economic-value profile, reconstructs the asset side from the same filing and
validates it against disclosures the model had not seen, and then recovers the
deposit assumption the published figure must have embedded. The result is not that
the assumption matters, which \citet{Chen2026} establishes, but that for this bank it
was both unrecoverable by any reader and, as Section~\ref{sec:sensitivity} shows
against the cross-bank evidence, implausible. The supervisory record indicates the
direction of the problem was recognised before it became acute \citep{FRB2023}; what
is added here is a measure of how much the disclosed number depended on something the
disclosure omitted.

A competing account of the 2023 failures should be noted, because this paper does not
test it. \citet{KellyRose2025} argue that the crisis is better explained by bank
business models, specifically exposure to the venture-capital and crypto-asset
sectors then under pressure, than by balance-sheet metrics alone, and that the
business-model account predicts which banks were affected more precisely. The
reconstruction here is a balance-sheet account. It is consistent with their critique
rather than a rebuttal of it: a measurement of what the filings could reveal says
nothing about which variable best separates the banks that failed from those that did
not.

Nothing here is a causal claim about the failure. The design establishes what was
measurable from published information. It does not test what any party would have done
with a 2022 economic-value figure, and the counterfactual in which SVB had Schwab's
funding mix is not observable.

\subsection{Structure}

The two balance sheets and the disclosure asymmetry between them come first, in
Section~\ref{sec:objects}, followed by the research questions and the
calibrate-then-freeze design in Section~\ref{sec:design} and the construction itself
in Section~\ref{sec:reconstruction}. Results arrive in two parts that should not be
confused: Section~\ref{sec:results} reports what the model produces, and
Section~\ref{sec:comparison} tests it against what the banks subsequently disclosed.
The deposit result sits in Section~\ref{sec:sensitivity}, which puts a number on every
assumption no filing constrains. Section~\ref{sec:interpretation} then separates what
the evidence establishes from what this paper reads into it, two biases against its
own argument included. Three appendices hold the curve diagnostics, the pricing
validation and the deposit solve.
